\documentclass[12pt,a4paper]{article}

\usepackage[utf8]{inputenc}
\usepackage[brazilian]{babel}
\usepackage{setspace}
\usepackage{mathptmx}
\usepackage{geometry}
\usepackage{url}
\usepackage{xcolor}
\usepackage{fancyhdr}
\usepackage{longtable}
\usepackage{enumitem}


\geometry{a4paper, left=2cm, right=2cm, bottom=4cm, top=3cm}
\setstretch{1.5}
\setlength{\parindent}{0pt}

\pagestyle{fancy}
\fancyhf{}

% ------------------------------------------------
% HEADER
% ------------------------------------------------
\lhead{Ata de Reunião}
\chead{\meetingcode}
\rhead{\meetingdate}

% ------------------------------------------------
% FOOTER
% ------------------------------------------------
\fancyfoot[L]{Confidencial}
\fancyfoot[C]{$\copyright$ \companyname, \the\year}
\fancyfoot[R]{Página \thepage}

\renewcommand{\footrulewidth}{0pt}

% ------------------------------------------------
% MACROS DA REUNIÃO
% ------------------------------------------------

\newcommand{\meetingtitle}{}
\newcommand{\meetingtype}{}
\newcommand{\meetingdate}{}
\newcommand{\meetingtime}{}
\newcommand{\meetinglocation}{}
\newcommand{\facilitator}{}
\newcommand{\notetaker}{}
\newcommand{\participants}{}
\newcommand{\meetingobjective}{}
\newcommand{\meetingcode}{}
\newcommand{\companyname}{}

\newcommand{\setmeetingtitle}[1]{\renewcommand{\meetingtitle}{#1}}
\newcommand{\setmeetingtype}[1]{\renewcommand{\meetingtype}{#1}}
\newcommand{\setmeetingdate}[1]{\renewcommand{\meetingdate}{#1}}
\newcommand{\setmeetingtime}[1]{\renewcommand{\meetingtime}{#1}}
\newcommand{\setmeetinglocation}[1]{\renewcommand{\meetinglocation}{#1}}
\newcommand{\setfacilitator}[1]{\renewcommand{\facilitator}{#1}}
\newcommand{\setnotetaker}[1]{\renewcommand{\notetaker}{#1}}
\newcommand{\setparticipants}[1]{\renewcommand{\participants}{#1}}
\newcommand{\setobjective}[1]{\renewcommand{\meetingobjective}{#1}}
\newcommand{\setmeetingcode}[1]{\renewcommand{\meetingcode}{#1}}
\newcommand{\setcompanyname}[1]{\renewcommand{\companyname}{#1}}


% ------------------------------------------------
% INÍCIO DO DOCUMENTO
% ------------------------------------------------

\begin{document}

\setmeetingcode{MIIA-ATA-f8c17192}

\setmeetingtitle{MIIA - Reunião de Planejamento da Sprint 1}
\setmeetingtype{SCRUM - Sprint Planning}
\setmeetingdate{11/02/2026}
\setmeetingtime{13:30 -- 14:30}
\setmeetinglocation{Sala de Reuniões do Laboratório de Modelagem e Simulação do IEAv}
\setfacilitator{Maj QOCon3 CCC \textbf{André} Francisco Morielo Caetano}
\setnotetaker{Maj QOCon3 CCC \textbf{André} Francisco Morielo Caetano}
\setcompanyname{Instituto de Estudos Avançados (IEAv)}

\setparticipants{
\begin{itemize}
\item Maj QOCon3 CCC \textbf{André} Francisco Morielo Caetano
\item Cap QOAV Lucas Silva \textbf{Lima}
\item 1\textsuperscript{o} Ten QOCon CMP R/2 Vinícius Guimarães Marques \textbf{Mattos}
\end{itemize}
}

\setobjective{Revisar os cards (atividades) previstas para Sprint, tendo em vista o objetivo maior (EPIC) proposto pela equipe: produzir um documento ``MIIA v1.0'' com design de alto e baixo nível como \emph{blueprint} para a arquitetura que iremos implementar depois}

% ------------------------------------------------
% TÍTULO
% ------------------------------------------------

\begin{center}
{\Large \textbf{\meetingtitle}}\\
\vspace{0.3em}
{\small Código da Ata: \meetingcode}
\end{center}

\vspace{1em}

\textbf{Tipo de Reunião:} \meetingtype

\textbf{Data:} \meetingdate

\textbf{Horário:} \meetingtime

\textbf{Local / Plataforma:} \meetinglocation

\textbf{Facilitador:} \facilitator

\textbf{Responsável pela Ata:} \notetaker

\vspace{1em}

\textbf{Participantes}

\participants

\vspace{1em}

\textbf{Objetivo da Reunião}

\meetingobjective

\vspace{1em}

\textbf{Pauta}

\begin{enumerate}
\item Avaliar cards com atividades propostas para a Sprint 1, associadas ao projeto MIIA.
\item Definir próximas ações e identificar os riscos associados às entregas propostas para a Sprint 1.
\end{enumerate}

\vspace{1em}

\textbf{Discussões}

\begin{itemize}
\item O Maj André conduziu a reunião apresentando inicialmente uma lista de cards sugeridos para a Sprint 1. Em seguida, cada card foi detalhado e avaliado em conjunto com os demais participantes, que puderam opinar e propor correções ou melhorias. Ao mesmo tempo, o Cap Lima trouxe um segundo conjunto de cards relacionados às atividades que ele vem conduzindo desde a Sprint 0 e que deverão prosseguir na Sprint 1; esses cartões também foram analisados coletivamente pelos presentes.
\end{itemize}

\vspace{1em}

\textbf{Decisões Tomadas}

\begin{itemize}
\item Os cards revisados foram considerados adequados pela equipe, embora hajam riscos (detalhados abaixo) que podem comprometer as entregas.
\item A equipe deve concentrar‑se nos entregáveis de maior valor na Sprint atual, mantendo os demais cartões no backlog para as próximas Sprints.
\end{itemize}

\vspace{1em}

\textbf{Ações Definidas}

\begin{longtable}{|p{6cm}|p{4cm}|p{3cm}|}
\hline
\textbf{Ação} & \textbf{Responsável} & \textbf{Prazo} \\
\hline
 Definir o Design de Alto Nível (HLD) do MIIA e apresentar para o Tech Lead & Maj QOCon3 CCC \textbf{André} & 13/02/2026 \\
\hline
 Trabalhar no PoC cliente/servidor do MIIA (gRPC/ONNX Runtime) & Cap QOAV \textbf{Lima} & 23/02/2026 \\
\hline
 Resolver, inicialmente, as pendências e atividades de Infraestrutura  & 1\textsuperscript{o} Ten QOCon CMP R/2 \textbf{Mattos} & 23/02/2026 \\
\hline
 Trabalhar nos demais cards  & Todos & 13/03/2026 \\
\hline
\end{longtable}

\vspace{1em}

\textbf{Bloqueios / Riscos}

\begin{itemize}
\item Maj André estará de férias de 18/02 a 27/02.
\item Cap Lima estará com dispensa por instalação de 23/02 a 04/03.
\item Ten Mattos estará sozinho de 23/02 a 27/02.
\end{itemize}

\vspace{2em}

\textbf{Encerramento}

Nada mais havendo a tratar, a reunião foi encerrada.

\vspace{3em}

\textbf{Assinaturas}

\vspace{1em}

\begin{longtable}{|p{6cm}|p{9cm}|}
\hline
\textbf{Nome} & \textbf{Assinatura} \\
\hline
Maj QOCon3 CCC \textbf{André} Francisco Morielo Caetano & \\[2.5cm]
\hline
Cap QOAV Lucas Silva \textbf{Lima} & \\[2.5cm]
\hline
1\textsuperscript{o} Ten QOCon CMP R/2 Vinícius Guimarães Marques \textbf{Mattos} & \\[2.5cm]
\hline
\end{longtable}

\end{document}